\documentclass[12pt,a4paper,oneside]{book}

% =====================================================================
% PACOTES ESSENCIAIS
% =====================================================================
\usepackage[brazilian]{babel}
\usepackage[T1]{fontenc}
\usepackage[utf8]{inputenc}
\usepackage{amsmath}
\usepackage{amssymb}
\usepackage{amsthm}
\usepackage{geometry}
\usepackage{fancyhdr}
\usepackage{hyperref}
\usepackage{booktabs}
\usepackage{mathtools}

% =====================================================================
% GEOMETRIA
% =====================================================================
\geometry{
    left=2.5cm,
    right=2.5cm,
    top=3cm,
    bottom=3cm,
    headheight=14pt
}

% =====================================================================
% CABECALHO E RODAPE
% =====================================================================
\pagestyle{fancy}
\fancyhf{}
\fancyhead[L]{\textit{Capitulo 3: Escala, Escalabilidade e Fractais}}
\fancyhead[R]{\thepage}
\fancyfoot[C]{\small Principios de Modelagem Matematica}
\renewcommand{\headrulewidth}{0.4pt}
\renewcommand{\footrulewidth}{0.4pt}

% =====================================================================
% AMBIENTES
% =====================================================================
\theoremstyle{definition}
\newtheorem{definition}{Definicao}[chapter]
\newtheorem{example}[definition]{Exemplo}
\newtheorem{remark}[definition]{Observacao}
\newtheorem{problem}{Problema}

% =====================================================================
% DOCUMENTO
% =====================================================================
\title{\textbf{Capitulo 3: Escala, Escalabilidade e Fractais}\\[0.4cm]
\large Principios de Modelagem Matematica}
\author{Baseado em Dym, C. L.}
\date{2026-04-14}

\begin{document}

\maketitle
\tableofcontents
\newpage

\chapter{Escala, Escalabilidade e Fractais}

\section{Ideia central do capitulo}

Neste capitulo, o foco e entender como um sistema muda quando alteramos seu tamanho caracteristico.
Esse tipo de pergunta aparece em engenharia, fisica, biologia e economia:

\begin{itemize}
    \item Se dobramos o tamanho, o efeito tambem dobra?
    \item Quais grandezas crescem mais rapido e se tornam dominantes?
    \item Existe uma lei de potencia do tipo $Y \propto X^\alpha$?
\end{itemize}

A analise de escala permite identificar termos relevantes sem resolver o problema completo.

\section{Escala e analise por ordem de grandeza}

\begin{definition}[Escala caracteristica]
Escala caracteristica e o valor tipico de uma variavel no problema, como comprimento $L$, velocidade $U$, tempo $T$ ou concentracao $C$.
\end{definition}

\begin{definition}[Ordem de grandeza]
Dizemos que duas quantidades sao da mesma ordem quando diferem por um fator pequeno (tipicamente menor que 10).
\end{definition}

\subsection{Nao dimensionalizacao rapida}

Considere a equacao de transporte 1D:
\begin{equation}
\frac{\partial c}{\partial t} + u\frac{\partial c}{\partial x} = D\frac{\partial^2 c}{\partial x^2}.
\end{equation}

Definimos variaveis adimensionais:
\begin{equation}
\hat{x} = \frac{x}{L}, \quad \hat{t} = \frac{t}{T}, \quad \hat{c} = \frac{c}{C_0}, \quad \hat{u} = \frac{u}{U}.
\end{equation}

Escolhendo $T=L/U$, obtemos:
\begin{equation}
\frac{\partial \hat{c}}{\partial \hat{t}} + \hat{u}\frac{\partial \hat{c}}{\partial \hat{x}} = \frac{1}{Pe}\frac{\partial^2 \hat{c}}{\partial \hat{x}^2},
\end{equation}
com numero de Peclet:
\begin{equation}
Pe = \frac{UL}{D}.
\end{equation}

\begin{remark}
Se $Pe \gg 1$, adveccao domina; se $Pe \ll 1$, difusao domina.
\end{remark}

\section{Leis de escala (power laws)}

Em muitos sistemas, observa-se:
\begin{equation}
Y = kX^\alpha,
\end{equation}
onde $\alpha$ e o expoente de escala.

Tomando logaritmo:
\begin{equation}
\log Y = \log k + \alpha \log X,
\end{equation}
logo, em grafico log-log, a inclinacao e $\alpha$.

\subsection{Exemplo 1: area e volume}

Para um corpo com comprimento caracteristico $L$:
\begin{equation}
A \propto L^2, \qquad V \propto L^3.
\end{equation}

Se aumentamos $L$ por fator 3:
\begin{equation}
A' = 3^2A = 9A, \qquad V' = 3^3V = 27V.
\end{equation}

\textbf{Conclusao}: volume cresce mais rapido que area. Isso explica, por exemplo, limites de troca termica em objetos grandes.

\subsection{Exemplo 2: periodo de pendulo}

Para pequenas amplitudes:
\begin{equation}
T = 2\pi\sqrt{\frac{\ell}{g}}.
\end{equation}

Logo:
\begin{equation}
T \propto \ell^{1/2}.
\end{equation}

Se $\ell$ quadruplica, o periodo dobra.

\section{Escalabilidade de modelos}

Escalabilidade significa como desempenho, custo ou erro variam com o tamanho do sistema.

\subsection{Exemplo 3: custo computacional}

Suponha um algoritmo com custo:
\begin{equation}
C(N) = aN^2.
\end{equation}

Se o tamanho de entrada dobra ($N \to 2N$):
\begin{equation}
C(2N) = a(2N)^2 = 4aN^2.
\end{equation}

Ou seja, dobrar o problema quadruplica o custo.

\subsection{Exemplo 4: transferencia de calor em objeto geometricamente similar}

Potencia termica aproximada por conveccao:
\begin{equation}
\dot{Q} \sim hA\Delta T.
\end{equation}

Como $A \propto L^2$, entao:
\begin{equation}
\dot{Q} \propto L^2.
\end{equation}

Energia armazenada (aprox. proporcional ao volume):
\begin{equation}
E \propto L^3.
\end{equation}

Tempo caracteristico de resfriamento:
\begin{equation}
\tau \sim \frac{E}{\dot{Q}} \propto L.
\end{equation}

Objetos maiores tendem a resfriar mais lentamente.

\section{Fractais e dimensao fractal}

Fractais sao estruturas com auto-similaridade: partes se parecem com o todo em diferentes escalas.

\begin{definition}[Dimensao fractal por auto-similaridade]
Se um objeto e composto por $N$ copias reduzidas por fator $r$ em cada escala,
a dimensao fractal e:
\begin{equation}
D = \frac{\log N}{\log(1/r)}.
\end{equation}
\end{definition}

\subsection{Exemplo 5: Curva de Koch}

Na Curva de Koch, cada segmento e substituido por $N=4$ segmentos com reducao $r=1/3$:
\begin{equation}
D = \frac{\log 4}{\log 3} \approx 1.2619.
\end{equation}

A curva tem dimensao maior que uma linha ($D=1$), mas menor que uma superficie ($D=2$).

\subsection{Exemplo 6: Triangulo de Sierpinski}

No Triangulo de Sierpinski: $N=3$ e $r=1/2$:
\begin{equation}
D = \frac{\log 3}{\log 2} \approx 1.585.
\end{equation}

\section{Roteiro pratico para estudar o capitulo 3}

\begin{enumerate}
    \item Identifique as variaveis e escalas caracteristicas do problema.
    \item Nao dimensionalize as equacoes.
    \item Extraia numeros adimensionais (Re, Pe, etc.).
    \item Compare ordens de grandeza para eliminar termos pequenos.
    \item Verifique se ha lei de potencia em grafico log-log.
    \item Quando houver auto-similaridade, estime dimensao fractal.
\end{enumerate}

\section{Exercicios curtos}

\begin{problem}
Se $Y \propto X^{1.7}$, por qual fator $Y$ muda quando $X$ dobra?
\end{problem}

\begin{problem}
Um objeto geometrico similar tem aresta multiplicada por 5. Como mudam area e volume?
\end{problem}

\begin{problem}
Em uma malha numerica 3D, o numero de celulas escala como $N \propto (L/\Delta x)^3$. Se $\Delta x$ cai pela metade, qual o fator de aumento em $N$?
\end{problem}

\section{Resumo final}

O capitulo 3 mostra que escalar nao e apenas aumentar tamanho: e entender como os mecanismos mudam de importancia.
A linguagem matematica dessa ideia aparece em:
\begin{itemize}
    \item leis de potencia,
    \item numeros adimensionais,
    \item eliminacao de termos por ordem de grandeza,
    \item geometria fractal e dimensao nao inteira.
\end{itemize}

Essas tecnicas tornam a modelagem mais robusta, economica e explicativa.

\end{document}
